% !TEX TS-program = lualatex
\documentclass[11pt, a4paper]{article}

% The shared preamble and bibliography live one directory up.
% Build from this directory: lualatex -shell-escape -halt-on-error jlt.tex
% Shared preamble of the papers in docs/paper (paper.tex, jlt/jlt.tex).
% Each paper adds its own \addbibresource (the path to references.bib is
% relative to the paper's directory).
% LuaLaTeX handles UTF-8 natively; fontspec lets the document pick up any
% Unicode glyphs present in the default font, including mathematical script
% characters used in listings (ℕ, ℤ, ℚ, ℝ, ℂ, Ø, and similar). We keep the
% `newunicodechar` fallbacks for body text so amsmath glyphs (\mathbb) are
% used in paragraphs.  Code listings use the `minted` package (Pygments-
% backed) with native UTF-8 — no per-character literate table needed.
\usepackage{fontspec}
\usepackage{newunicodechar}
\usepackage[english]{babel}    % Language support
\usepackage{geometry}          % Adjust page margins
\geometry{margin=1in}

% --- Recommended Math & Science Packages ---
\usepackage{amsmath, amssymb, amsthm} % Essential math tools
\usepackage{mathtools}
\usepackage{mathrsfs}

% Theorem environments (amsthm is loaded above).
\newtheorem{theorem}{Theorem}
\newtheorem{definition}{Definition}
\newtheorem{corollary}{Corollary}
\newtheorem{lemma}{Lemma}
\newtheorem{proposition}{Proposition}
\newtheorem{conjecture}{Conjecture}
\usepackage{graphicx}                  % Support for images
\usepackage{hyperref}                  % Clickable links & cross-refs
\usepackage{booktabs}                  % High-quality tables
% csquotes is loaded later — after `minted` (which pulls in `fvextra`)
% to silence the fvextra/csquotes load-order warning.

\newunicodechar{∩}{\ensuremath{\cap}}
\newunicodechar{ℚ}{\ensuremath{\mathbb{Q}}}
\newunicodechar{ℝ}{\ensuremath{\mathbb{R}}}
\newunicodechar{𝔽}{\ensuremath{\mathbb{F}}}
\newunicodechar{∃}{\ensuremath{\exists}}
\newunicodechar{π}{\ensuremath{\pi}}
\newunicodechar{⊂}{\ensuremath{\subset}}
\newunicodechar{⊆}{\ensuremath{\subseteq}}
\newunicodechar{∈}{\ensuremath{\in}}
\newunicodechar{∉}{\ensuremath{\notin}}
\newunicodechar{∧}{\ensuremath{\land}}
\newunicodechar{∨}{\ensuremath{\lor}}
\newunicodechar{≤}{\ensuremath{\leq}}
\newunicodechar{≥}{\ensuremath{\geq}}
\newunicodechar{μ}{\ensuremath{\mu}}
\newunicodechar{α}{\ensuremath{\alpha}}
\newunicodechar{β}{\ensuremath{\beta}}
\newunicodechar{λ}{\ensuremath{\lambda}}
\newunicodechar{Ω}{\ensuremath{\Omega}}
\newunicodechar{𝔸}{\ensuremath{\mathbb{A}}}
\newunicodechar{↔}{\ensuremath{\leftrightarrow}}
\newunicodechar{χ}{\ensuremath{\chi}}
\newunicodechar{Σ}{\ensuremath{\Sigma}}
\newunicodechar{Τ}{\ensuremath{\mathrm{T}}}

% For code listings
\usepackage{xcolor}
% Landscape tables:
\usepackage[figuresright]{rotating}
\usepackage[table]{xcolor}
\usepackage{pdflscape}
\usepackage{tabularx}
\usepackage{booktabs}


% Graphs
\usepackage{graphviz}

\usepackage{tikz}
\usetikzlibrary{shapes.geometric, arrows.meta, positioning, calc, backgrounds}
\usepackage{tikz-cd}

\usepackage{draftwatermark}
\SetWatermarkScale{5}

\usepackage[style=numeric, sorting=none]{biblatex}

% `minted` (Pygments-backed) handles UTF-8 natively in code listings —
% no per-character literate table needed.  Requires `lualatex -shell-escape`
% (the build pipeline in docs/paper/Makefile and the CI install step in
% the project Makefile are updated accordingly).
\usepackage{minted}
\usepackage{csquotes}  % Load after minted/fvextra to silence load-order warning.

% Use JuliaMono for code listings — designed specifically for Unicode-rich
% source (Julia uses the full math repertoire) with native coverage of
% Greek (φ, Σ, Τ, χ, ...), math double-struck (ℕ, ℤ, ℚ, ℝ, ℂ, 𝔻, 𝔽, ...),
% math-fraktur, math-script, sub/superscripts, and arrows.  Replaces the
% earlier `lmmono10` (no Greek) and `Fira Mono` (no math-double-struck)
% choices made in the lstlisting era.  Installed locally via
% `brew install --cask font-juliamono`; CI install adds `fonts-juliamono`
% (Ubuntu 22.04+) under ci-install-paper-deps in the project Makefile.
% Scale=0.90 keeps glyph metrics close to the document's serif body.
\setmonofont{JuliaMono}[Scale=0.90]

\setminted{
  fontsize=\small,
  breaklines=true,
  breakanywhere=true,
  frame=single,
  framesep=2mm,
  autogobble,
}
\usemintedstyle{friendly}

% C++ is the default code-listing language in this paper; declare short
% block + inline aliases so source markup reads `\begin{cpp}...\end{cpp}`
% and `\cppinline{...}` rather than the verbose `\begin{minted}{cpp}` /
% `\mintinline{cpp}{...}` forms.  The optional `[name]` first argument
% names the environment directly; without it minted would produce
% `\begin{cppcode}` and `\cpp{...}`, which would collide with our \cpp
% short alias usage.
\newminted[cpp]{cpp}{}
\newmintinline[cppinline]{cpp}{}

\definecolor{codegreen}{rgb}{0,0.6,0}
\definecolor{codegray}{rgb}{0.5,0.5,0.5}
\definecolor{codepurple}{rgb}{0.58,0,0.82}

\newcommand{\bind}{\mathbin{>\!\!>\mkern-6.7mu=}}
\newcommand{\fish}{\mathbin{\gg}}
\newcommand{\Cpp}[1]{C\kern-.05em\raise.3ex\hbox{\footnotesize ++}#1}

\addbibresource{../references.bib}
\graphicspath{{../figures/}}

\title{The Juliet Posture: Structurally Typed Functional Programming in Standard \Cpp{23}}
\author{Vincent Kr\"autler}
\date{\today}

\begin{document}
\maketitle

\begin{abstract}
Parametric polymorphism, type classes, structural recursion and algebraic laws are usually associated with languages built for them: Haskell, OCaml, Agda.
We show that they are available in standard \Cpp{23}, with a stock compiler and no extensions.
The vehicle is a fragment of the language we call the \emph{Juliet Posture} ($\mathbf{Jlt}$), defined by three rules: values are \cppinline{std::semiregular} and functions \cppinline{constexpr} and pure; generic code varies with a type only through that type's concepts and traits; and recursion only folds, over a term, a finite carrier or explicit fuel.
Under these rules a \Cpp{} program translates into System~F: concepts and traits by the dictionary-passing translation of type classes, recursion by the fold encoding of data types, so programs in the fragment are parametric and terminate.
A side effect is that one program runs at compile time or at runtime, unchanged.
As a running example, algebraic laws of functions (injectivity, monotonicity, involution) are stated once as \cppinline{constexpr} values and read twice: by concepts at compile time, by an ordinary fold at runtime.
The two readers are the two halves of the dictionary-passing translation, side by side.
The rules are a discipline the compiler does not enforce; we state where the correspondence holds, where it is a translation rather than an isomorphism, and what the fragment gives up.
\end{abstract}

% Section skeleton (#981).  Each section names what it must establish; prose
% follows once the calculus (appendix) is drafted.

\section{Introduction}
\label{sec:intro}
% Establishes: the claim (fancy FP with a recent clang++), the audience (C++
% first, System F as the warranty), the contributions, the roadmap.

\section{The Posture}
\label{sec:posture}
% Establishes: the three rules, each with a small listing, and what each buys
% (referential transparency; parametricity via dictionaries; termination via
% folds).  Compliance is at the level of the paper's listings.

\section{The Translation}
\label{sec:translation}
% Establishes: what a reader must believe about the System F reading, by one
% example and the rules it uses.  The full calculus (Appendix~\ref{app:fjlt})
% is the warranty; a reader who stops here knows the claims and their cost.

\subsection{One shape, two monoids}
\label{sec:translation:example}

Here is the whole translation in one shape.
A monoid is a carrier with an associative operation and a unit for it.
Listing~\ref{lst:monoid} states that shape as a concept over a carrier \cppinline{T} and an operation \cppinline{Op}, and then shows that \cppinline{std::uint64_t} has it twice.

% Exhibit test pending: identity_v / total_v registrations for the transparent
% std::plus<> and friends (the tree keys them on std::plus<T>).
\begin{listing}[ht]
\begin{cpp}
template <typename T, typename Op>
concept Monoid = std::semiregular<T> && Total<T, Op>
              && requires(const T a, const T b) {
  { Op{}(a, b) }        -> std::same_as<T>;
  { identity_v<T, Op> } -> std::same_as<T>;
};

static_assert(Monoid<std::uint64_t, std::plus<>>);
static_assert(Monoid<std::uint64_t, std::multiplies<>>);

template <typename T, typename Op = std::plus<>>
  requires Monoid<T, Op>
constexpr T twice(T x) { return Op{}(x, x); }

static_assert(twice<std::uint64_t>(21) == 42);
static_assert(twice<std::uint64_t, std::multiplies<>>(21) == 441);
\end{cpp}
\begin{minted}{text}
MonoidDict T Op  =  { total : 𝟙,  unit : T,  op : T → T → T }
plus_u64   =  { total = ∗,  unit = 0,  op = λa. λb. a + b }
times_u64  =  { total = ∗,  unit = 1,  op = λa. λb. a · b }
twice      =  Λ T. Λ Op. λ d : MonoidDict T Op. λ x : T.  d.op x x
twice [u64] [plus]  plus_u64  21
twice [u64] [times] times_u64 21
\end{minted}
\caption{The monoid shape, and two monoids on one carrier; below it, the same program in System~F, line for line.
The shape has four requirements of three kinds: a standard concept (rule~1), a declared judgment (\cppinline{Total}, a trait), an operation the compiler reads off the carrier, and a unit, which is a choice and therefore also a trait.
The operation is part of the index, so one carrier has an additive and a multiplicative monoid, and nobody wrote a struct for either.}
\label{lst:monoid}
\end{listing}

The lower half of the listing reads the upper half as a System~F program.
The concept is a record type, one field per requirement.
A \cppinline{requires} clause on a function is one more parameter, the record.
A call passes the record, and the compiler finds the record to pass.
That is dictionary passing, the translation of type classes into System~F,\footnote{Wadler and Blott~\cite{wadler1989adhoc}; the record is their dictionary.} and Table~\ref{tab:dictionary} names the \Cpp{} devices it meets.
Two of them deserve a sentence each.
\cppinline{Op} is a type with exactly one value, the closure \cppinline{Op{}}, so under translation the type \emph{is} that closure, and \cppinline{Op{}(a, b)} is \texttt{d.op a b}.
And a unit is a choice.
The compiler can see that \cppinline{std::uint64_t} has a \cppinline{+}, but not which element is neutral for it, so someone states the fact, once, as a trait specialisation.
In Haskell that someone writes an \texttt{instance}; here they write
\begin{cpp}
template <> inline constexpr std::uint64_t identity_v<std::uint64_t, std::plus<>> = 0;
\end{cpp}
and the one-definition rule makes sure they write it once.

\begin{table}[ht]
\centering
\small
\begin{tabularx}{\linewidth}{@{}XXX@{}}
\toprule
\Cpp{23} & type-class reading & under translation \\
\midrule
concept & class & record type \\
requirement met by the type's own operation & an instance nobody wrote & field filled by a primitive \\
trait specialisation & instance & field supplied by a value \\
constrained trait specialisation & instance with a context & function from records to a record \\
\cppinline{requires} clause on a function & class constraint & one more parameter \\
\bottomrule
\end{tabularx}
\caption{Dictionary passing, met in \Cpp{}.
The second row has no Haskell counterpart: concepts are structural, so a type that has the operation satisfies the shape without a declaration.}
\label{tab:dictionary}
\end{table}

The operation being part of the index is a small thing with a known payoff.
Haskell indexes an instance by the carrier alone, so one carrier has one monoid and the second needs a wrapper type (\texttt{Sum}, \texttt{Product}).
Here \cppinline{Monoid<T, Op>} is a relation between a carrier and an operation, and the two \cppinline{static_assert}s above are its two witnesses at one carrier.
The defaulted \cppinline{Op} picks the additive reading when a caller does not care.

\subsection{Shapes compose by conjunction}
\label{sec:translation:semiring}

A semiring is two monoids on one carrier, with laws tying them together.
Listing~\ref{lst:semiring} states it as a conjunction of shapes and tests it on two carriers, accepting both, and tests the monoid shape on two more, rejecting both.

\begin{listing}[ht]
\begin{cpp}
template <typename T, typename Add, typename Mul>
concept Semiring = CommutativeMonoid<T, Add> && Monoid<T, Mul>
                && Distributive<T, Mul, Add> && Annihilating<T, Mul, Add>;

static_assert( Semiring<std::uint64_t, std::plus<>, std::multiplies<>>);
static_assert( Semiring<bool, std::logical_or<>, std::logical_and<>>);
static_assert(!Monoid<bool, std::plus<>>);
static_assert(!Monoid<int, std::plus<>>);
\end{cpp}
\caption{The semiring shape as a conjunction, and four verdicts.
$(\mathbb{Z}/2^{64}, +, \cdot)$ and $(\mathbb{B}, \lor, \land)$ are semirings.
$(\mathbb{B}, +)$ is not a monoid, because \cppinline{true + true} is the \cppinline{int} 2 and the return-type requirement fails.
$(\mathbb{Z}, +)$ on \cppinline{int} is not a monoid, because signed overflow is undefined behaviour, so \cppinline{Total<int, std::plus<>>} is declared false.}
\label{lst:semiring}
\end{listing}

The four verdicts fill a $2 \times 2$ table, and the table is the point of the listing.

\begin{center}
\begin{tabular}{@{}lll@{}}
\toprule
 & accepted & rejected \\
\midrule
structurally, read off the type & $(\mathbb{Z}/2^{64}, +, \cdot)$, $(\mathbb{B}, \lor, \land)$ & $(\mathbb{B}, +)$: not closed \\
by declaration, a trait records a judgment & $0$ and $1$ as the units & $(\mathbb{Z}, +)$ on \cppinline{int}: not total \\
\bottomrule
\end{tabular}
\end{center}

Structural where the type has the operation, declared where a fact has to be recorded.
Both kinds are requirements of the same shape, and both become fields of the same record.
Under translation the conjunction is a product, \texttt{SemiringDict T Add Mul = MonoidDict T Add × MonoidDict T Mul × \ldots}, so a conjunct that is itself a shape becomes a field holding a record, which is how a superclass is passed.

One thing the listing does not do, and cannot.
A concept checks signatures.
It sees that \cppinline{Add} and \cppinline{Mul} exist on \cppinline{T} with the right types, and that a unit is declared for each.
It cannot see that \cppinline{Mul} distributes over \cppinline{Add}, and the \cppinline{bool} row shows why that matters: $(\mathbb{B}, \lor, \land)$ distributes, $(\mathbb{B}, +, \cdot)$ is not even closed, and the shape only notices the second.
\cppinline{Distributive<T, Mul, Add>} is therefore not a fact the compiler decides about \cppinline{T}.
It is a fact someone states, like the unit, and Section~\ref{sec:laws} is about where such facts come from, who reads them, and what can be inferred from them.
The shape permits \cppinline{identity_v<T, Add> == identity_v<T, Mul>}, as a semiring should; the one-element carrier witnesses it.

\subsection{Instances with contexts}
\label{sec:translation:contexts}
% A constrained trait specialisation: the constant maps into a monoid form a
% monoid pointwise, and the unit of the composite needs the unit of the
% codomain (identity_registry<ConstantMorphism<A, B>, Op> requires
% IsPointed<B, Op>).  A function from dictionaries to a dictionary.  The
% bridge to Section~\ref{sec:laws}: the readers on a composite are such
% functions.

\subsection{The rules the argument uses}
\label{sec:translation:rules}
% Only the rules the example exercised: concept -> record type; trait
% specialisation -> dictionary (value, or function of its context); requires
% clause -> parameter; qualified call -> application; fold with fuel -> Church
% eliminator.  The remaining rules are mechanical and sit in
% Appendix~\ref{app:fjlt:translation}.

\subsection{What the translation guarantees}
\label{sec:translation:theorems}
% The three theorem STATEMENTS (type preservation, simulation, coherence as the
% ODR), with the corollaries (termination, parametricity).  Proof sketches:
% Appendix~\ref{app:fjlt:theorems}.

\subsection{From \Cpp{} to the calculus}
\label{sec:translation:adequacy}
% The adequacy argument in prose: which C++23 features the fragment keeps,
% which it forbids, and the gap (the compiler does not enforce the rules).

\section{Laws of Functions}
\label{sec:laws}
% Establishes: the running example.  Laws as constexpr values (a profile); the
% declared facts (closure, the sign table, entailment); the formers (pairing,
% copairing; product and sum by their defining forms); the two readers as the
% two halves of the dictionary translation.

\section{Two Binding Times}
\label{sec:binding}
% Establishes: the same program at compile time and at runtime; Unknown as
% the price of structural recursion on fuel.

\section{Limits and Seams}
\label{sec:limits}
% Establishes: the rules are a discipline, not enforced; a translation, not an
% isomorphism; the adequacy argument and its gap; the seams found while
% building the example (#981).

\section{Related Work}
\label{sec:related}
% Establishes: F^G (Siek, Lumsdaine) as the nearest neighbour; the
% Featherweight line; type classes and concepts; formal semantics of C++
% fragments.  What is new.

\section{Conclusion}
\label{sec:conclusion}

\appendix
% Featherweight Jlt (FJlt): the core calculus of the Jlt fragment, its
% translation into System F, and the three theorems with proof sketches
% (#981, option B; full proofs and mechanisation: #849).

\section{Featherweight Jlt}
\label{app:fjlt}
% Design (decided on #981): no overlapping instances; fold and explicit fuel as
% the only recursion; laws as ordinary values; multi-parameter instances with
% contexts.

\subsection{Syntax}
\label{app:fjlt:syntax}

\subsection{Typing}
\label{app:fjlt:typing}

\subsection{Evaluation}
\label{app:fjlt:evaluation}

\subsection{Translation into System F}
\label{app:fjlt:translation}

\subsection{Theorems}
\label{app:fjlt:theorems}
% Type preservation; simulation; coherence.  Corollaries: termination
% (Girard), parametricity (Reynolds).


\printbibliography
\end{document}
